\documentclass{IOS-Book-Article}

% Define \runtitle if not provided by class
\providecommand{\runtitle}[1]{}

\usepackage{mathptmx}
\usepackage[utf8]{inputenc}
\usepackage[T1]{fontenc}
\usepackage[vietnamese,english]{babel}

% Math and symbols
\usepackage{amsmath,amssymb}
\usepackage{amsthm}
\newtheorem{definition}{Definition}

% Tables and figures
\usepackage{booktabs}
\usepackage{multirow}
\usepackage{graphicx}
\usepackage{float}
\usepackage{enumitem}

% TikZ for diagrams
\usepackage{tikz}
\usetikzlibrary{shapes.geometric, arrows.meta, positioning, fit, backgrounds, calc}

% pgfplots for charts
\usepackage{pgfplots}
\pgfplotsset{compat=1.18}

% Tighten spacing for page limit (no titlesec to avoid cls conflict)
\makeatletter
\renewcommand\section{\@startsection{section}{1}{\z@}%
  {-1.5ex \@plus -0.5ex \@minus -.3ex}%
  {0.8ex \@plus .2ex}%
  {\normalfont\large\bfseries}}
\renewcommand\subsection{\@startsection{subsection}{2}{\z@}%
  {-1.2ex \@plus -0.3ex \@minus -.2ex}%
  {0.5ex \@plus .1ex}%
  {\normalfont\normalsize\bfseries}}
\renewcommand\subsubsection{\@startsection{subsubsection}{3}{\z@}%
  {-1.0ex \@plus -0.2ex \@minus -.2ex}%
  {0.4ex \@plus .1ex}%
  {\normalfont\normalsize\itshape}}
\makeatother
\setlength{\textfloatsep}{8pt plus 2pt minus 2pt}
\setlength{\floatsep}{6pt plus 2pt minus 2pt}
\setlength{\intextsep}{6pt plus 2pt minus 2pt}
\setlength{\abovecaptionskip}{4pt}
\setlength{\belowcaptionskip}{2pt}

\begin{document}

\pagestyle{headings}
\def\thepage{}

\begin{frontmatter}

\title{A Three-Tier Knowledge Graph RAG for Vietnamese Legal Question Answering}

% TODO: Add ORCID iDs for each author
\author[A]{\fnms{Hong Hien} \snm{Le}%
\thanks{Corresponding Author: Le Hong Hien, Faculty of Computer Science, University of Information Technology, Vietnam National University Ho Chi Minh City, Vietnam; E-mail: hienlh.16@grad.uit.edu.vn.}},
\author[A]{\fnms{Dinh Hien} \snm{Nguyen}},
\author[B]{\fnms{Quoc Hung} \snm{Ngo}}
and
\author[A]{\fnms{Viet Dung} \snm{Dang}}

\runningauthor{H.H. Le et al.}
\runtitle{Three-Tier KG RAG for Vietnamese Legal QA}
\address[A]{Faculty of Computer Science, University of Information Technology, Vietnam National University Ho Chi Minh City, Vietnam}
\address[B]{School of Computer Science, University College Dublin, Ireland}

\begin{abstract}
To answer legal questions accurately, a system must navigate hierarchical document structures, follow cross-references, and reason across multiple hops, but standard Retrieval Augmented Generation (RAG) systems cannot do this. Vietnamese legal text makes this gap worse. Tonal language, domain-specific abbreviations, and a rigid structural hierarchy (Law$\to$Chapter$\to$Article$\to$Clause$\to$Point) all compound retrieval difficulty. We propose a three-tier knowledge graph RAG architecture that extends the Legal-Onto Rela-model with a formal knowledge model $\mathcal{M} = (\mathcal{O}, \mathcal{G}, \mathcal{T}, \Phi)$, integrating an OWL ontology for terminological knowledge, a domain-specific knowledge graph for assertional knowledge, and a document tree forest for structural navigation. The three retrieval tiers, specifically LLM-guided tree traversal, dual-level semantic retrieval with intent-aware Personalized PageRank, and Reciprocal Rank Fusion, operate over these representations. To evaluate the method, we conducted experiments on 200 legal questions in Vietnam belonging to two different legal fields: enterprise law (150 questions on company formation, capital, dissolution) and traffic law (50 questions on violations, penalties, procedures) across 13 legal documents. Our system achieves 84.5\% Hit@10 (MRR 0.461), outperforming TF-IDF (40.5\%) and BM25 (36.0\%) by substantial margins. Per-domain results confirm generalization: 83.3\% Hit@10 on enterprise law and 88.0\% on traffic law, showing the flexibility and strong querying capabilities across various legal domains. We are not aware of prior work combining ontology-enhanced knowledge graphs with multi-tier retrieval for Vietnamese legal question answering.
\end{abstract}

\begin{keyword}
Knowledge Graph\sep Legal Question Answering\sep Retrieval Augmented Generation\sep Ontology\sep Vietnamese NLP\sep Multi-Tier Retrieval\sep Multi-Domain
\end{keyword}

\end{frontmatter}

\section{Introduction}

Legal forums and social networks in Vietnam are flooded with questions on business formation, traffic penalties, and regulatory compliance. Yet answers are often fragmented, contradictory, or cite repealed provisions.

Legal documents exhibit \textbf{complex hierarchical structures}: a single law organizes content into parts, chapters, sections, articles, clauses, and points, where each level carries distinct legal scope. They also contain pervasive \textbf{cross-references}: answering ``Can a foreign investor establish a limited liability company in Vietnam?'' requires following reference chains across multiple articles to assemble conditions, procedures, and exceptions.

The \textbf{Vietnamese legal domain} adds further complexity. Vietnamese is tonal with ambiguous compound word boundaries. Legal practitioners use extensive abbreviations (\textit{HDQT} for ``Hoi dong quan tri'', \textit{CTCP} for ``Cong ty co phan'') that general models miss. Vietnamese law follows a rigid hierarchy such as Constitution, Codes, Decrees, and Circulars, where document type determines legal authority.

Traditional RAG systems \cite{r1} retrieve passages by vector similarity between query and document embeddings \cite{r2}. This works adequately for many domains \cite{r3,r4}, but legal QA exposes its weaknesses. Legal terminology carries precise meanings that general embeddings blur (\textbf{semantic gap}); documents organize content hierarchically, yet vector search treats them as flat text (\textbf{no structural awareness}); and answering a question often requires chaining references across articles---something dense retrieval cannot do (\textbf{no multi-hop reasoning}).

Existing approaches tackle parts of this problem. LightRAG \cite{r5} and GraphRAG \cite{r6} model entity relationships but disregard document hierarchy. PageIndex \cite{r7} exploits structure but has no semantic reasoning component. Legal ontology systems \cite{r8,r9} formalize domain knowledge but remain disconnected from modern retrieval pipelines. We propose a \textbf{three-tier knowledge graph RAG} architecture that brings these capabilities together:
\begin{enumerate}
\item \textbf{LLM-guided tree traversal} through document hierarchies for structural navigation (Tier~1),
\item \textbf{Dual-level semantic retrieval} with intent-aware Personalized PageRank on a domain knowledge graph (Tier~2),
\item \textbf{Reciprocal Rank Fusion} for cross-tier merging with knowledge graph expansion (Tier~3).
\end{enumerate}

We contribute a three-tier RAG architecture that combines structural, semantic, and graph-based retrieval. We also build a Vietnamese legal KG with 5,936 entities (14 types) and 5,633 relations (56 types) using LLM extraction and Vietnamese slug normalization. Our multi-intent PPR algorithm adapts edge weights to query intents. To validate domain flexibility, we evaluate across two legal domains: enterprise law and traffic safety law, using 13 legal documents. Experiments on 200 real questions show 84.5\% Hit@10 (+44.0pp over the best baseline), with consistent performance across both domains.

\section{Related Work}
\label{sec:related}

\textbf{RAG and KG-augmented retrieval.}
RAG \cite{r1} grounds generation in retrieved evidence. Since early dense retrieval pipelines \cite{r10}, designs have grown more sophisticated: Self-RAG \cite{r11} adds reflection, CRAG \cite{r12} corrective retrieval, and HyDE \cite{r13} hypothetical document generation; surveys \cite{r3,r4} trace this progression. On the knowledge graph side, GraphRAG \cite{r6} enables multi-hop reasoning and LightRAG \cite{r5} offers lighter entity-relation extraction, but both treat documents as bags of entities and lose structural context. Embedding-based KG methods (TransE \cite{r14}, RotatE \cite{r15}) capture relational patterns yet are limited to single-hop prediction. Personalized PageRank \cite{r16,r17} computes query-specific node importance. PageIndex \cite{r7} takes a different path with LLM-guided tree traversal---achieving high precision with few articles, but missing content outside the traversal path.

\textbf{Legal NLP.}
Legal IR has moved from lexical methods \cite{r18,r19} toward domain-adapted transformers such as LEGAL-BERT \cite{r20}, which capture legal phrasing better than general models but still treat documents as flat passages. Ontology-based approaches \cite{r9} add richer structure; Legal-Onto \cite{r8} combines the Rela-model \cite{r21} with conceptual graphs for Vietnamese Land Law, though it relies on manual inference rules and was not designed for retrieval. For Vietnamese specifically, PhoBERT \cite{r22} and VnCoreNLP \cite{r23} provide foundational NLP capabilities, and Ngo et al.\ \cite{r24} combined BM25 with BERT \cite{r25} and domain heuristics---an effective baseline that nonetheless lacks graph-based reasoning. Surveys \cite{r26,r27} confirm that bridging formal legal knowledge with scalable retrieval remains an open problem.

\textbf{Positioning.}
Table~\ref{tab:related-comparison} maps each system's capabilities. The gap motivating our work is the absence of a system that combines ontological knowledge representation with multi-tier retrieval for Vietnamese legal text.

\begin{table}[H]
\centering\footnotesize
\caption{Feature comparison of legal QA approaches. $^*$Not a RAG system; included for ontology comparison.}
\label{tab:related-comparison}
\resizebox{\columnwidth}{!}{%
\begin{tabular}{llcccccc}
\toprule
\textbf{Category} & \textbf{Feature} & \textbf{NaiveRAG} & \textbf{LightRAG} & \textbf{GraphRAG} & \textbf{PageIndex} & \textbf{Legal-Onto$^*$} & \textbf{Ours} \\
\midrule
\multirow{3}{*}{Retrieval} & Semantic Retrieval & \checkmark & \checkmark & -- & -- & -- & \checkmark \\
& Graph Traversal & -- & \checkmark & \checkmark & -- & -- & \checkmark \\
& Tree Navigation & -- & -- & -- & \checkmark & -- & \checkmark \\
\midrule
\multirow{2}{*}{Knowledge} & Knowledge Graph & -- & \checkmark & \checkmark & -- & \checkmark & \checkmark \\
& Domain Ontology & -- & -- & -- & -- & \checkmark & \checkmark \\
\midrule
\multirow{2}{*}{Reasoning} & Multi-hop & -- & -- & \checkmark & -- & \checkmark & \checkmark \\
& Intent-aware & -- & -- & -- & -- & -- & \checkmark \\
\midrule
\multirow{3}{*}{Integration} & Multi-tier & -- & -- & -- & -- & -- & \checkmark \\
& Rank Fusion & -- & -- & -- & -- & -- & \checkmark \\
& Multi-domain & -- & -- & -- & -- & -- & \checkmark \\
\bottomrule
\end{tabular}}%
\end{table}

\section{Methodology}
\label{sec:method}

\subsection{Problem Formulation}

Given a legal corpus $\mathcal{D} = \{d_1, \ldots, d_m\}$ with hierarchically organized articles and a query $q$, retrieve a ranked list $A^* = [a_1^*, \ldots, a_k^*]$ of relevant articles. We define $A^* = \Phi(q, \mathcal{M})$ where $\mathcal{M} = (\mathcal{O}, \mathcal{G}, \mathcal{T}, \Phi)$ integrates ontology $\mathcal{O}$ (terminological), knowledge graph $\mathcal{G}$ (assertional), document forest $\mathcal{T}$ (structural), and retrieval functions $\Phi$.

\subsection{System Overview}

Figure~\ref{fig:architecture} illustrates the architecture. The \textbf{offline phase} constructs $\mathcal{M}$ through parsing, extraction, and ontology generation. The \textbf{online phase} processes queries via ontology expansion (Stage~0) followed by three retrieval tiers leveraging $\mathcal{T}$, $\mathcal{G}$, and $\mathcal{O}$.

% Figure 1: System Architecture
\begin{figure}[H]
\centering
\resizebox{0.85\textwidth}{!}{%
\begin{tikzpicture}[
    % Styles
    box/.style={rectangle, draw, rounded corners, minimum height=0.7cm, minimum width=1.8cm, align=center, font=\small},
    phase/.style={rectangle, draw, dashed, rounded corners, inner sep=8pt},
    tier/.style={rectangle, draw, dashed, rounded corners, inner sep=6pt, fill=blue!3},
    arrow/.style={-{Stealth[length=2mm]}, thick},
    dataarrow/.style={-{Stealth[length=2mm]}, thick, blue!60},
    label/.style={font=\footnotesize\bfseries},
]

% === OFFLINE PHASE (Left) ===
\node[box, fill=orange!20] (doc) at (0,0) {Legal\\Document};
\node[box, fill=blue!15] (parse) at (3,0) {Document\\Parsing};
\node[box, fill=blue!15] (extract) at (6,0) {LLM Entity \&\\Relation Extraction};
\node[box, fill=green!15] (kgbuild) at (3,-2.5) {Incremental KG\\Builder (Dedup)};
\node[box, fill=green!15] (summary) at (6,-2.5) {Chapter \&\\Article Summaries};
\node[box, fill=purple!20, minimum width=3cm] (kg) at (4.5,-5) {Knowledge Model $\mathcal{M}$\\$(\mathcal{O}, \mathcal{G}, \mathcal{T}, \Phi)$};

% Offline arrows
\draw[arrow] (doc) -- (parse);
\draw[arrow] (parse) -- (extract);
\draw[arrow] (extract.south) -- ++(0,-0.5) -| (kgbuild.north);
\draw[arrow] (extract.south) -- ++(0,-0.5) -| (summary.north);
\draw[arrow] (kgbuild.south) -- (kg.north west);
\draw[arrow] (summary.south) -- (kg.north east);

% Offline phase box
\begin{scope}[on background layer]
\node[phase, fill=gray!5, fit=(doc)(parse)(extract)(kgbuild)(summary)(kg), label={[label]above:Offline Phase: Knowledge Model Construction}] (offline) {};
\end{scope}

% === ONLINE PHASE (Right) ===
\node[box, fill=yellow!20] (query) at (10.5,0) {User\\Query};
\node[box, fill=blue!15] (analyzer) at (13.5,0) {Query Analyzer\\(Intent + Type)};

% Tier 1
\node[box, fill=green!15] (tree) at (11,-2.5) {Tree Traversal\\(LLM-guided)};

% Tier 2
\node[box, fill=green!15] (dual) at (16,-2.5) {Dual-Level\\Retriever};
\node[box, fill=green!15] (ppr) at (16,-4.0) {Intent-Aware\\PPR};

% Tier 3
\node[box, fill=red!15] (bridge) at (13.5,-5.8) {Semantic Bridge\\(RRF Fusion)};

% Output
\node[box, fill=purple!20] (gen) at (13.5,-7.5) {Answer\\Generation};
\node[box, fill=yellow!20] (answer) at (16.5,-7.5) {Legal\\Answer};

% Online arrows
\draw[arrow] (query) -- (analyzer);
\draw[arrow] (analyzer.south) -- ++(0,-0.5) -| (tree.north);
\draw[arrow] (analyzer.south) -- ++(0,-0.5) -| (dual.north);
\draw[arrow] (dual) -- (ppr);
\draw[arrow] (tree.south) |- (bridge.west);
\draw[arrow] (ppr.south) |- (bridge.east);
\draw[arrow] (bridge) -- (gen);
\draw[arrow] (gen) -- (answer);

% KG to Online connections
\coordinate (splitpt) at (8.5,-5);
\draw[blue!60, thick] (kg.east) -- (splitpt);
\draw[blue!60, thick] (splitpt) -- (8.5,-2.5);
\draw[dataarrow] (8.5,-2.5) -- (tree.west);
\draw[dataarrow] (splitpt) -- (8.5,-5.8) -- (bridge.west);

% Tier labels
\node[font=\scriptsize\itshape, gray, anchor=west] at (12.1,-2.5) {Tier 1};
\node[font=\scriptsize\itshape, gray, anchor=east] at (14.9,-2.5) {Tier 2};
\node[font=\scriptsize\itshape, gray] at (13.5,-5.1) {Tier 3};

% Online phase box
\begin{scope}[on background layer]
\node[phase, fill=gray!5, fit=(query)(analyzer)(tree)(dual)(ppr)(bridge)(gen)(answer), label={[label]above:Online Phase: Three-Tier Query Processing}] (online) {};
\end{scope}

\end{tikzpicture}
}
\caption{System architecture. Offline: knowledge model construction. Online: three-tier query processing (Tier~1: tree traversal, Tier~2: semantic+PPR, Tier~3: RRF fusion).}
\label{fig:architecture}
\end{figure}

\subsection{Knowledge Representation Model}

Our model extends Legal-Onto \cite{r8}, which defines $K = (C, R, \textit{RULES}) + (\textit{Key}, \textit{Rel}, \textit{weight})$ integrating the Rela-model \cite{r21} with conceptual graphs. We extend this by separating terminological from assertional knowledge, adding structural representation, and replacing manual rules with LLM-based retrieval.

\begin{definition}[Legal Knowledge Model]
Given a legal corpus $\mathcal{D}$, the knowledge model is:
\begin{equation}
\mathcal{M} = (\mathcal{O}, \mathcal{G}, \mathcal{T}, \Phi)
\end{equation}
where $\mathcal{O}$ is a domain ontology, $\mathcal{G}$ is a knowledge graph, $\mathcal{T}$ is a document forest, and $\Phi$ is a set of retrieval functions.
\end{definition}

\textbf{Ontology} $\mathcal{O} = (C, P, H, L)$ represents terminological knowledge: $C$~=~OWL classes for legal entity types, $P = P_\text{obj} \cup P_\text{data}$~=~object and datatype properties, $H \subseteq C \times C$~=~class hierarchy via \texttt{rdfs:subClassOf} (4-level taxonomy), and $L: C \cup P \to \text{Vietnamese}$~=~label mapping (e.g., \texttt{JointStockCompany} $\mapsto$ ``cong ty co phan''). This separates the TBox from Legal-Onto's flat concept set, enabling class hierarchy traversal for query expansion.

\textbf{Knowledge Graph} $\mathcal{G} = (E, R)$ represents assertional knowledge. Each entity $e = (\textit{id}, \textit{name}, \textit{type}, \textit{conf}, \textit{src})$ has a slug-normalized ID, one of 14 types (Organization, Person/Role, Legal Term, Legal Reference, Condition, Action, Penalty, Time Period, Monetary, etc.), confidence score, and source provenance. Each relation $r = (e_i, \textit{pred}, e_j, \textit{evid}, \textit{conf})$ carries evidence text and confidence across 56 types with formal properties: symmetric $\{\textit{related\_to}, \textit{contradicts}, \textit{synonym}\}$ and transitive $\{\textit{is\_a}, \textit{part\_of}, \textit{contains}\}$. Compared to Legal-Onto's conceptual graph with binary similarity vectors, $\mathcal{G}$ provides richer typing, confidence tracking, and evidence-backed relations.

\textbf{Example: ontology class vs.\ KG entity.} Consider ``joint-stock company'' (\textit{cong ty co phan}). In $\mathcal{O}$, this is an OWL \textbf{class} \texttt{JointStockCompany} with a taxonomy path \texttt{JointStockCompany} $\xrightarrow{\text{subClassOf}}$ \texttt{Company} $\xrightarrow{\text{subClassOf}}$ \texttt{LegalEntity}; querying ``joint-stock company'' traverses $H$ to also match sibling classes (\texttt{LLC}, \texttt{Partnership}), enabling \textit{query expansion}. In $\mathcal{G}$, the same concept is an \textbf{entity node} \textit{cong-ty-co-phan} (type: Organization) linked to legal facts: \textit{requires}~$\to$~Art.\,111 (establishment conditions), \textit{has\_authority}~$\to$~\textit{hoi-dong-quan-tri} (board of directors). PPR traverses these edges to retrieve relevant articles. In short, $\mathcal{O}$ answers ``\textit{what kind of entity is this?}'' while $\mathcal{G}$ answers ``\textit{what legal rules apply to it?}''. This distinction has practical impact: for the query ``Does a joint-stock company need a power-of-attorney when establishing an LLC subsidiary?'', $\mathcal{O}$ recognizes both \texttt{JointStockCompany} and \texttt{LLC\_SingleMember} as sibling subclasses of \texttt{Company}, expanding the query to match articles about both types; $\mathcal{G}$ then links entity \textit{nguoi-dai-dien-phan-von-gop} (authorized representative) via \textit{governed\_by} to Art.\,79 (delegation rules). No baseline retrieves any relevant article for this query (Hit@30\,=\,0), while our system ranks the correct article first (MRR\,=\,1.0).

\textbf{Document Forest} $\mathcal{T} = \{T_1, \ldots, T_m\} \cup X$ represents structural knowledge. Each $T_i$ is a rooted tree with six levels (Document${}\to{}$\allowbreak Chapter${}\to{}$\allowbreak Section${}\to{}$\allowbreak Article${}\to{}$\allowbreak Clause${}\to{}$\allowbreak Point), and $X$ contains cross-reference edges between documents. The forest spans multiple legal domains, with documents organized into domain groups (e.g., enterprise law, traffic safety law) that enable domain-aware routing during retrieval.

\textbf{Retrieval Functions} $\Phi = \{\phi_0, \phi_1, \phi_2, \phi_3\}$ replace Legal-Onto's manual inference rules:
\begin{itemize}[nosep,leftmargin=1.5em]
\item $\phi_0\colon Q \times \mathcal{O} \to Q'$ (ontology expansion),
\item $\phi_1\colon Q' \times \mathcal{T} \to 2^A$ (tree traversal),
\item $\phi_2\colon Q' \times \mathcal{G} \times \mathcal{O} \to 2^A$ (semantic retrieval),
\item $\phi_3\colon 2^A \times 2^A \times \mathcal{G} \to A^*$ (fusion).
\end{itemize}
Where Legal-Onto applies hand-crafted deductive rules, our functions leverage LLM reasoning while preserving formal relation properties in the graph structure.

\subsection{Knowledge Graph Construction}

The offline phase builds $\mathcal{G}$ through four steps. \textbf{(1)~Document Parsing:} Legal documents are parsed into hierarchical trees following Vietnamese formatting standards, forming the document forest $\mathcal{T}$ with position, parent-child relationships, and cross-references. \textbf{(2)~Entity-Relation Extraction:} A single LLM call per article extracts entities (14 types) and relations simultaneously, producing evidence-grounded triples with source citations. \textbf{(3)~Incremental KG Building:} Entities are deduplicated via Vietnamese slug normalization: Unicode NFD decomposition, diacritics removal, and lowercasing (e.g., ``Cong ty co phan'' $\to$ \textit{cong-ty-co-phan}), which uses a 4-pass pipeline: exact cache lookup, synonym matching (17 entries), abbreviation expansion (25 entries), and new slug creation. Cross-document entity deduplication ensures consistent entity resolution across all 13 documents. \textbf{(4)~Summary Generation:} LLM-generated chapter and article summaries enable Tier~1's tree traversal decisions.

\subsubsection{Intent-Aware Edge Weights}

Table~\ref{tab:intent-weights} defines intent-specific weight multipliers for PPR.

\begin{table}[H]
\centering\footnotesize
\caption{Intent-relation weight matrix.}
\label{tab:intent-weights}
\begin{tabular}{lccccc}
\toprule
\textbf{Relation} & \textbf{Definition} & \textbf{Procedure} & \textbf{Right} & \textbf{Obligation} & \textbf{Penalty} \\
\midrule
defines\_as & 1.0 & 0.5 & 0.5 & 0.5 & 0.5 \\
includes & 0.95 & 0.5 & 0.5 & 0.5 & 0.5 \\
requires & 0.5 & 1.0 & 0.5 & 0.85 & 0.5 \\
has\_authority & 0.5 & 0.5 & 1.0 & 0.5 & 0.5 \\
has\_obligation & 0.5 & 0.5 & 0.5 & 1.0 & 0.5 \\
has\_penalty & 0.5 & 0.5 & 0.5 & 0.5 & 1.0 \\
\bottomrule
\end{tabular}
\end{table}

This allows PPR to traverse paths most relevant to query intent (e.g., \textit{defines\_as} for definitions, \textit{requires} for procedures).

\subsection{Three-Tier Retrieval}

\subsubsection{Stage 0: Query Analysis}

The query analyzer performs: (1)~\textbf{intent detection} across 6 types (Penalty, Definition, Procedure, Requirement, Reference, General) via Vietnamese regex with LLM fallback; (2)~\textbf{type classification} into 7 categories; (3)~\textbf{abbreviation expansion} (e.g., \textit{TNHH} $\to$ \textit{trach nhiem huu han}); and (4)~\textbf{ontology expansion} $\phi_0(q, \mathcal{O})$ via class hierarchy traversal in $H$ with Vietnamese label lookup through $L$.

\subsubsection{Tier 1: LLM-Guided Tree Traversal}

Tier~1 navigates $\mathcal{T}$ using LLM reasoning in three loops: \textbf{Loop~0} (Forest$\to$Domain$\to$Document) first matches the query to a domain group using pre-computed domain keywords, then selects top~5 documents within the matched domain; \textbf{Loop~1} (Document$\to$Chapter) selects top~8 chapters using summaries and keywords; \textbf{Loop~2} (Chapter$\to$Article) identifies top~7 articles from article summaries. General Provisions chapters are automatically included. Each loop produces a reasoning path for interpretable retrieval.

\subsubsection{Tier 2: Dual-Level Semantic Retrieval}

Tier~2 performs hybrid semantic search over $\mathcal{G}$ and $\mathcal{O}$ with six components. \textit{Low-level} (entity-specific): (1)~keyphrase matching, (2)~query-entity cosine similarity, (3)~intent-aware PPR, (4)~ontology class matching. \textit{High-level} (theme/concept): (5)~legal theme identification, (6)~hierarchy traversal through $\mathcal{O}$'s taxonomy. The combined score is:
\begin{equation}
S_{\text{T2}} = .05 S_{\text{kp}} + .20 S_{\text{con}} + .25 S_{\text{PPR}} + .20 S_{\text{sem}} + .15 S_{\text{thm}} + .15 S_{\text{hier}}
\end{equation}

\textbf{Multi-Intent Personalized PageRank.} The PPR component computes query-specific node importance on $\mathcal{G}$:
\begin{equation}
\text{PPR}(v) = (1-\alpha) \sum_{u \in N(v)} \frac{w(u,v) \cdot \text{PPR}(u)}{d(u)} + \alpha \cdot p(v)
\end{equation}
where $\alpha=0.15$ is the teleport probability, $w(u,v)$ is the edge weight adjusted by query intent (Table~\ref{tab:intent-weights}), $d(u)$ is the weighted out-degree of $u$, and $p(v)$ is the personalization vector seeded from matched entities.

\subsubsection{Tier 3: Semantic Bridge (RRF Fusion)}

Tier~3 merges results from Tier~1 and Tier~2 using Reciprocal Rank Fusion \cite{r28}:
\begin{equation}
\text{RRF}(d) = \sum_{s \in \{T1, T2, KG\}} w_s \cdot \frac{1}{k + \text{rank}_s(d)}
\end{equation}
where $k=60$ and source weights are $w_{T1}=1.4$, $w_{T2}=0.7$, $w_{KG}=1.0$. Tree traversal receives the highest weight as it provides the most reliable structural signal for legal documents, while KG expansion supplements with cross-reference articles that tree traversal may miss. RRF handles heterogeneous score distributions (LLM discrete selections vs.\ continuous semantic scores) without calibration.

\textbf{Knowledge graph expansion.} After initial fusion, we expand results through $\mathcal{G}$'s relation edges. Strong relations (\textit{references}, \textit{requires}, \textit{depends\_on}) are prioritized for cross-chapter linking (up to 5 articles), while expansion relations (\textit{defined\_as}, \textit{related\_to}) handle same-chapter augmentation (up to 2 articles). Cross-document references are explicitly resolved via document-prefixed entity IDs.

\textbf{Adjacent article expansion.} After RRF scoring, we expand $\pm 2$ articles from each result with decay factor 0.7, as correct answers often lie in neighboring articles. Top-ranked articles are passed to an LLM for answer generation.

\section{Experiments}
\label{sec:experiments}

\subsection{Experimental Setup}

\textbf{Dataset.} We evaluate on 13 Vietnamese legal documents spanning two legal domains: (1)~\textit{Enterprise law}: Vietnam Enterprise Law 2020 (No.~59/2020/QH14) and 9 related decrees covering business registration, corporate governance, and dissolution; (2)~\textit{Traffic safety law}: Traffic Safety Law 2024 (No.~36/2024/QH15) and 2 penalty decrees covering traffic rules, violations, and administrative procedures. Including 77 chapters and a total of 840 articles. The evaluation set: 200 questions from Vietnamese legal platforms (LuatVietnam, ThuVienPhapLuat), expert-annotated (1--8 relevant articles per question, avg.\ 2.3), comprising 150 enterprise law and 50 traffic safety law questions.

\textbf{Knowledge graph statistics.} Table~\ref{tab:kg-stats} summarizes the constructed knowledge graph.

\begin{table}[H]
\centering\footnotesize
\caption{Knowledge graph statistics.}
\label{tab:kg-stats}
\begin{tabular}{lr@{\hskip 1.2cm}lr@{\hskip 1.2cm}lr}
\toprule
\textbf{Entity Type} & \textbf{Count} & \textbf{Relation Type} & \textbf{Count} & \textbf{Structure} & \textbf{Count} \\
\midrule
Total Entities & 5,936 & Total Relations & 5,633 & Document Tree Nodes & 9,226 \\
Action & 1,903 & requires & 1,377 & Documents & 13 \\
Legal Term & 1,632 & applies\_to & 810 & Chapter Summaries & 77 \\
Organization & 500 & has\_authority & 506 & Article Summaries & 840 \\
Legal Reference & 451 & references & 434 & Cross-references & 168 \\
Person/Role & 392 & condition\_for & 361 & & \\
Penalty & 327 & includes & 358 & & \\
Other (8 types) & 731 & Other (50 types) & 1,787 & & \\
\bottomrule
\end{tabular}
\end{table}

\textbf{Implementation.} Claude~3.5 Haiku for tree traversal, query analysis and entity extraction. Embeddings: \texttt{paraphrase-multilingual-MiniLM-L12-v2} \cite{r2}. PPR: iterative computation, max~50 iterations, tolerance~$10^{-6}$. All LLM responses cached in SQLite ($\sim$2.1\,GB) for reproducibility.

\textbf{Metrics.} Hit@$k$, Recall@$k$, Precision@$k$, NDCG@$k$ at $k \in \{1,3,5,10,20,30\}$, and MRR.

\subsection{Baselines}

\textbf{BM25} \cite{r18} ($k_1\!=\!1.2$, $b\!=\!0.75$); \textbf{TF-IDF} \cite{r19} with Vietnamese segmentation; \textbf{Semantic Search} \cite{r2} using multilingual sentence-transformers; \textbf{LightRAG} \cite{r5} with default parameters; \textbf{PageIndex} \cite{r7} (our Tier~1 in isolation). All baselines retrieve from the full 13-document corpus with max\_results=30.

\subsection{Main Results}

Table~\ref{tab:main-results} presents Hit@$k$ performance.

\begin{table}[H]
\centering\footnotesize
\caption{Hit@$k$ performance (\%). Best \textbf{bold}, 2nd \underline{underlined}. -- = pending evaluation.}
\label{tab:main-results}
\begin{tabular}{l@{\hskip 3pt}r@{\hskip 3pt}r@{\hskip 3pt}r@{\hskip 3pt}r@{\hskip 3pt}r}
\toprule
\textbf{Method} & \textbf{Hit@5} & \textbf{Hit@10} & \textbf{Hit@20} & \textbf{Hit@30} & \textbf{MRR} \\
\midrule
Semantic & 10.0 & 15.5 & 24.0 & 25.5 & 0.144 \\
Keyword & 22.0 & 33.0 & 51.5 & 64.0 & 0.161 \\
BM25 & 21.0 & 36.0 & 58.5 & 67.5 & 0.144 \\
TF-IDF & \underline{26.0} & \underline{40.5} & \underline{59.5} & \underline{71.0} & \underline{0.191} \\
LightRAG & -- & -- & -- & -- & -- \\
PageIndex & -- & -- & -- & -- & -- \\
\textbf{3-Tier (ours)} & \textbf{65.0} & \textbf{84.5} & \textbf{87.5} & \textbf{90.5} & \textbf{0.461} \\
\bottomrule
\end{tabular}
\end{table}

Semantic search scores only 15.5\% Hit@10 due to general embeddings failing on Vietnamese legal semantics. BM25/TF-IDF achieve moderate performance ($\sim$36--41\%) with weak ranking. Our 3-Tier system reaches 84.5\% Hit@10 (+44.0pp over TF-IDF, the best available baseline), demonstrating that structural navigation combined with graph-based semantic retrieval produces dramatically better results than flat retrieval methods.

Figure~\ref{fig:hitk-comparison} visualizes the consistent advantage across all $k$ values.

\begin{figure}[H]
\centering
\begin{tikzpicture}
\begin{axis}[
    width=\columnwidth,
    height=4cm,
    ybar,
    bar width=3.5pt,
    ylabel={Hit Rate (\%)},
    symbolic x coords={@5,@10,@20,@30},
    xtick=data,
    ymin=0, ymax=100,
    enlarge x limits=0.2,
    legend style={at={(0.5,-0.3)}, anchor=north, legend columns=3, font=\scriptsize},
    every axis plot/.append style={fill opacity=0.85},
    grid=major,
    grid style={dashed, gray!30},
]
\addplot[fill=gray!40, draw=gray!60] coordinates {(@5,10.0) (@10,15.5) (@20,24.0) (@30,25.5)};
\addplot[fill=green!30, draw=green!50] coordinates {(@5,22.0) (@10,33.0) (@20,51.5) (@30,64.0)};
\addplot[fill=blue!30, draw=blue!50] coordinates {(@5,21.0) (@10,36.0) (@20,58.5) (@30,67.5)};
\addplot[fill=cyan!30, draw=cyan!50] coordinates {(@5,26.0) (@10,40.5) (@20,59.5) (@30,71.0)};
\addplot[fill=red!50, draw=red!70] coordinates {(@5,65.0) (@10,84.5) (@20,87.5) (@30,90.5)};
\legend{Semantic, Keyword, BM25, TF-IDF, 3-Tier (ours)}
\end{axis}
\end{tikzpicture}
\caption{Hit@$k$ comparison across all methods.}
\label{fig:hitk-comparison}
\end{figure}

\subsection{Detailed IR Metrics}

Table~\ref{tab:detailed-metrics} reports detailed metrics. MRR is 0.461 (first relevant article at avg.\ rank~$\sim$2.2). Recall reaches 62.0\% at $k\!=\!30$.

\begin{table}[H]
\centering\footnotesize
\caption{Detailed IR metrics for 3-Tier system.}
\label{tab:detailed-metrics}
\begin{tabular}{lrrrr}
\toprule
\textbf{$k$} & \textbf{Hit} & \textbf{Recall} & \textbf{Precision} & \textbf{NDCG} \\
\midrule
1 & 31.5 & 12.1 & 31.5 & 31.5 \\
3 & 50.5 & 29.1 & 25.8 & 30.3 \\
5 & 65.0 & 39.6 & 21.1 & 34.6 \\
10 & 84.5 & 55.2 & 14.6 & 41.3 \\
20 & 87.5 & 59.4 & 7.9 & 42.8 \\
30 & 90.5 & 62.0 & 5.5 & 43.5 \\
\bottomrule
\end{tabular}
\end{table}

\subsection{Multi-Domain Performance}

A key design goal is domain flexibility: the same architecture should generalize to new legal domains without retraining or manual rule adjustment. Table~\ref{tab:domain-results} compares retrieval performance across the two legal domains.

\begin{table}[H]
\centering\footnotesize
\caption{Per-domain Hit@$k$ (\%) for 3-Tier system.}
\label{tab:domain-results}
\begin{tabular}{lrrrrrrr}
\toprule
\textbf{Domain} & \textbf{$n$} & \textbf{Hit@1} & \textbf{Hit@5} & \textbf{Hit@10} & \textbf{Hit@20} & \textbf{Hit@30} & \textbf{MRR} \\
\midrule
Enterprise Law & 150 & 29.3 & 62.7 & 83.3 & 86.7 & 88.7 & 0.436 \\
Traffic Safety Law & 50 & 38.0 & 72.0 & 88.0 & 90.0 & 96.0 & 0.535 \\
\midrule
\textbf{Combined} & \textbf{200} & \textbf{31.5} & \textbf{65.0} & \textbf{84.5} & \textbf{87.5} & \textbf{90.5} & \textbf{0.461} \\
\bottomrule
\end{tabular}
\end{table}

Traffic safety law achieves higher Hit@10 (88.0\% vs.\ 83.3\%) and MRR (0.535 vs.\ 0.436) than enterprise law. This is attributable to the more structured nature of traffic penalty articles, which map clearly to specific violations with explicit penalty amounts and procedures. Enterprise law queries, particularly those involving cross-decree references and multi-step procedures, pose greater retrieval challenges. Crucially, no domain-specific tuning was applied, the same KG construction pipeline, retrieval parameters, and RRF weights serve both domains, confirm the generality of the architecture.

\subsection{Analysis}

\subsubsection{Performance by Query Category}

Table~\ref{tab:category-results} breaks down Hit@10 by question category.

\begin{table}[H]
\centering\footnotesize
\caption{Hit@10 by question category (\%).}
\label{tab:category-results}
\begin{tabular}{lrr}
\toprule
\textbf{Category} & \textbf{Hit@10} & \textbf{Count} \\
\midrule
\multicolumn{3}{l}{\textit{Enterprise Law}} \\
Capital contributions & 100.0 & 18 \\
Share/capital transfer & 100.0 & 11 \\
Dissolution \& bankruptcy & 95.7 & 47 \\
Legal representatives & 83.3 & 6 \\
Business registration changes & 78.6 & 28 \\
Company formation & 69.2 & 13 \\
Company types & 69.2 & 13 \\
Other enterprise & 42.9 & 14 \\
\midrule
\multicolumn{3}{l}{\textit{Traffic Safety Law}} \\
Traffic administrative procedures & 92.0 & 25 \\
Traffic violation penalties & 91.7 & 12 \\
Traffic rules & 75.0 & 8 \\
Other traffic & 80.0 & 5 \\
\bottomrule
\end{tabular}
\end{table}

Capital-related queries (100\%) benefit most from Tier~1's direct navigation to specific articles. Dissolution queries (95.7\%) are well-served by the structured penalty decree hierarchy. Traffic violation penalties (91.7\%) demonstrate PPR's effectiveness: penalty queries activate the \textit{has\_penalty} intent weight, guiding graph traversal toward sanction-related articles. Categories with lower performance (company formation 69.2\%, company types 69.2\%) involve multi-decree cross-references where the relevant article resides in a decree outside the primary law.

\subsubsection{Qualitative and Error Analysis}

\textbf{Success (multi-hop):} ``If a member borrows money to meet charter capital and later withdraws it, what are the penalties?'' (Tier~1 navigates to Article~47 on member obligations; Tier~2's PPR follows \textit{requires} edges from the ``charter capital'' entity to Article~16\S5 listing prohibited acts including fraudulent capital declaration; Tier~3 fuses both in top~5, achieving \textbf{Hit@5}).
\textbf{Success (cross-chapter):} ``A partnership member contributed a house worth 1\,billion VND; it is now worth 1.5\,billion---can they withdraw it?'' (Tier~1 reaches Articles~178--180 on partnership rules; Tier~2 adds Article~186 on capital transfer via \textit{condition\_for} edges, correctly identifying that withdrawal is prohibited but transfer is permitted; \textbf{Hit@5}).
\textbf{Failure:} ``Requirements for foreign-invested enterprise capital structure?'' (requires synthesizing Investment Law articles not in our corpus; Tier~2 partially reaches related articles but ranks at position~15).

We manually categorized the 31 failed queries (15.5\%). The dominant cause is \textbf{cross-corpus gaps} (12 queries, 38.7\%): the answer resides in a law or decree outside the 13-document corpus (e.g., Investment Law, Tax Law). \textbf{Compound term splitting} accounts for 9 failures (29.0\%): Vietnamese word segmentation fragments multi-syllable legal terms, degrading both embedding quality and entity matching. The remaining 10 failures (32.3\%) stem from \textbf{ambiguous intent}: queries mixing multiple intent types cause PPR to favor one traversal path over another.

\textbf{Reproducibility.} All LLM responses are cached in SQLite ($\sim$2.1\,GB), making results deterministic (std.\ dev.\ = 0).

\section{Discussion}
\label{sec:discussion}

The multi-domain results (Table~\ref{tab:domain-results}) demonstrate that the three-tier architecture generalizes across legal domains without modification. Traffic safety law (88.0\% Hit@10) and enterprise law (83.3\%) both achieve strong performance using identical retrieval parameters, KG construction pipelines, and RRF weights. This is significant because the two domains differ substantially in structure: enterprise law involves complex cross-decree references and multi-step procedures, while traffic law features well-structured penalty tables with explicit mappings between violations and sanctions. The architecture handles both patterns because Tier~1 provides domain-appropriate structural navigation while Tier~2's intent-aware PPR adapts its traversal to query semantics.

Intent-aware PPR proves particularly effective for penalty and procedure queries. By weighting edges according to query intent (Table~\ref{tab:intent-weights}), the algorithm surfaces paths that generic graph traversal would rank lower---a ``procedure'' query and a ``penalty'' query about the same entity need different neighborhoods of the knowledge graph. Traffic law penalty queries (91.7\% Hit@10) illustrate this well: the \textit{has\_penalty} weight guides PPR directly to sanction articles.

The RRF weight configuration ($w_{T1}\!=\!1.4$, $w_{T2}\!=\!0.7$, $w_{KG}\!=\!1.0$) reflects an empirical finding: tree traversal is the most reliable signal for Vietnamese legal documents, whose rigid hierarchy means that correct structural navigation almost always reaches the relevant chapter. Dual-level semantic retrieval serves as a complementary signal, catching articles reachable only via cross-references or thematic relations that lie outside the tree's traversal path.

The knowledge model $\mathcal{M}$, separating terminological ($\mathcal{O}$), assertional ($\mathcal{G}$), and structural ($\mathcal{T}$) representations, provides the scaffolding for integrating both approaches. This separation also facilitates multi-domain deployment: adding a new legal domain requires only parsing new documents and running the same KG construction pipeline---no architectural changes or retraining.

\section{Conclusion}
\label{sec:conclusion}

This paper introduced a three-tier knowledge graph RAG system for Vietnamese legal QA, built on a formal knowledge model $\mathcal{M} = (\mathcal{O}, \mathcal{G}, \mathcal{T}, \Phi)$ that extends Legal-Onto by separating terminological, assertional, and structural representations. On 200 real questions spanning two legal domains, the system reaches 84.5\% Hit@10 (MRR 0.461), improving over the strongest baseline by 44.0pp. The central finding is that structural navigation and semantic graph retrieval are complementary, and this complementarity holds across domains: 83.3\% Hit@10 on enterprise law and 88.0\% on traffic safety law, with no domain-specific tuning.

\textbf{Limitations.} Tier~1 tree traversal requires multiple LLM calls, adding $\sim$7\,s of latency per query. Entity extraction has not been validated against a gold-standard annotation. The two domains tested share Vietnamese legal formatting conventions; generalization to structurally different legal traditions remains untested.

\textbf{Future work.} We plan to extend the knowledge graph to Investment, Tax, and Labor law to further validate multi-domain generalization. We also aim to add multi-turn conversational support, investigate reasoning chains beyond two hops, and produce explainable retrieval traces that cite the graph paths used in ranking.

% References ordered by first citation appearance (Vancouver/NLM style)
\begin{thebibliography}{28}
\setlength{\itemsep}{0pt}\setlength{\parsep}{0pt}\setlength{\parskip}{0pt}

\bibitem{r1}
Lewis P, Perez E, Piktus A, et al. Retrieval-Augmented Generation for Knowledge-Intensive NLP Tasks. In: Advances in Neural Information Processing Systems 33 (NeurIPS 2020); 2020. p. 9459-74.

\bibitem{r2}
Reimers N, Gurevych I. Sentence-BERT: Sentence Embeddings using Siamese BERT-Networks. In: Proceedings of the 2019 Conference on Empirical Methods in Natural Language Processing (EMNLP); 2019 Nov 3-7; Hong Kong. 2019. p. 3982-92, doi: 10.18653/v1/D19-1410.

\bibitem{r3}
Gao Y, Xiong Y, Gao X, et al. Retrieval-Augmented Generation for Large Language Models: A Survey. arXiv preprint. 2024. arXiv:2312.10997, doi: 10.48550/arXiv.2312.10997.

\bibitem{r4}
Fan W, Ding Y, Ning L, et al. A Survey on RAG Meeting LLMs: Towards Retrieval-Augmented Large Language Models. In: Proceedings of the 30th ACM SIGKDD Conference on Knowledge Discovery and Data Mining (KDD 2024); 2024. p. 6491-501, doi: 10.1145/3637528.3671470.

\bibitem{r5}
Guo Z, et al. LightRAG: Simple and Fast Retrieval-Augmented Generation. arXiv preprint. 2024. arXiv:2410.05779, doi: 10.48550/arXiv.2410.05779.

\bibitem{r6}
Edge D, Trinh H, Cheng N, et al. From Local to Global: A Graph RAG Approach to Query-Focused Summarization. arXiv preprint. 2024. arXiv:2404.16130, doi: 10.48550/arXiv.2404.16130.

\bibitem{r7}
Zhang M, Tang Y, et al. PageIndex: Next-Generation Vectorless, Reasoning-based RAG. 2025. Available from: https://pageindex.ai/blog/pageindex-intro.

\bibitem{r8}
Nguyen T, Nguyen H, Pham V, Tran D, Selamat A. Legal-Onto: An Ontology-based Model for Representing the Knowledge of a Legal Document. In: Proceedings of the 17th International Conference on Evaluation of Novel Approaches to Software Engineering (ENASE 2022); 2022. p. 426-34, doi: 10.5220/0011088400003176.

\bibitem{r9}
Sartor G, et al. Approaches to Legal Ontologies: Theories, Domains, Methodologies. Dordrecht: Springer; 2011. 272 p., doi: 10.1007/978-94-007-0120-5.

\bibitem{r10}
Karpukhin V, Oguz B, Min S, et al. Dense Passage Retrieval for Open-Domain Question Answering. In: Proceedings of the 2020 Conference on Empirical Methods in Natural Language Processing (EMNLP); 2020. p. 6769-81, doi: 10.18653/v1/2020.emnlp-main.550.

\bibitem{r11}
Asai A, Wu Z, Wang Y, Sil A, Hajishirzi H. Self-RAG: Learning to Retrieve, Generate, and Critique through Self-Reflection. In: Proceedings of the International Conference on Learning Representations (ICLR 2024); 2024, doi: 10.48550/arXiv.2310.11511.

\bibitem{r12}
Yan SQ, et al. Corrective Retrieval Augmented Generation. arXiv preprint. 2024. arXiv:2401.15884, doi: 10.48550/arXiv.2401.15884.

\bibitem{r13}
Gao L, Ma X, Lin J, Callan J. Precise Zero-Shot Dense Retrieval without Relevance Labels. In: Proceedings of the 61st Annual Meeting of the Association for Computational Linguistics (ACL 2023); 2023. p. 1762-77, doi: 10.18653/v1/2023.acl-long.99.

\bibitem{r14}
Bordes A, Usunier N, Garcia-Duran A, Weston J, Yakhnenko O. Translating Embeddings for Modeling Multi-relational Data. In: Advances in Neural Information Processing Systems 26 (NeurIPS 2013); 2013. p. 2787-95.

\bibitem{r15}
Sun Z, Deng Z-H, Nie J-Y, Tang J. RotatE: Knowledge Graph Embedding by Relational Rotation in Complex Space. In: Proceedings of the International Conference on Learning Representations (ICLR 2019); 2019, doi: 10.48550/arXiv.1902.10197.

\bibitem{r16}
Gleich D. PageRank Beyond the Web. SIAM Rev. 2015;57(3):321-63, doi: 10.1137/140976649.

\bibitem{r17}
Jeh G, Widom J. Scaling Personalized Web Search. In: Proceedings of the 12th International Conference on World Wide Web (WWW 2003); 2003 May 20-24; Budapest, Hungary. New York: ACM; 2003. p. 271-9, doi: 10.1145/775152.775191.

\bibitem{r18}
Robertson S, Zaragoza H. The Probabilistic Relevance Framework: BM25 and Beyond. Found Trends Inf Retr. 2009;3(4):333-89, doi: 10.1561/1500000019.

\bibitem{r19}
Manning C, Raghavan P, Schutze H. Introduction to Information Retrieval. New York: Cambridge University Press; 2008. 544 p., doi: 10.1017/CBO9780511809071.

\bibitem{r20}
Chalkidis I, Fergadiotis M, Malakasiotis P, Aletras N, Androutsopoulos I. LEGAL-BERT: The Muppets Straight out of Law School. In: Findings of the Association for Computational Linguistics: EMNLP 2020; 2020. p. 2898-904, doi: 10.18653/v1/2020.findings-emnlp.261.

\bibitem{r21}
Do N, Nguyen H, Selamat A. Knowledge-Based model of Expert Systems using Rela-model. Int J Softw Eng Knowl Eng. 2018;28(8):1047-90, doi: 10.1142/S0218194018500304.

\bibitem{r22}
Nguyen DQ, Nguyen AT. PhoBERT: Pre-trained Language Models for Vietnamese. In: Findings of the Association for Computational Linguistics: EMNLP 2020; 2020. p. 1037-42, doi: 10.18653/v1/2020.findings-emnlp.92.

\bibitem{r23}
Vu T, Nguyen DQ, Nguyen DQ, Dras M, Johnson M. VnCoreNLP: A Vietnamese Natural Language Processing Toolkit. In: Proceedings of the 2018 Conference of the North American Chapter of the Association for Computational Linguistics: Demonstrations (NAACL 2018); 2018. p. 56-60, doi: 10.18653/v1/N18-5012.

\bibitem{r24}
Ngo H, Nguyen T, Nguyen D, Pham M. AimeLaw at ALQAC 2021: Enriching Neural Network Models with Legal-Domain Knowledge. In: Proceedings of the 13th International Conference on Knowledge and Systems Engineering (KSE 2021). IEEE; 2021, doi: 10.1109/KSE53942.2021.9648712.

\bibitem{r25}
Devlin J, Chang M, Lee K, Toutanova K. BERT: Pre-training of Deep Bidirectional Transformers for Language Understanding. In: Proceedings of the 2019 Conference of the North American Chapter of the Association for Computational Linguistics: Human Language Technologies (NAACL-HLT 2019); 2019. p. 4171-86, doi: 10.18653/v1/N19-1423.

\bibitem{r26}
Zhong H, Xiao C, Tu C, Zhang T, Liu Z, Sun M. How Does NLP Benefit Legal System: A Summary of Legal Artificial Intelligence. In: Proceedings of the 58th Annual Meeting of the Association for Computational Linguistics (ACL 2020); 2020. p. 5218-30, doi: 10.18653/v1/2020.acl-main.466.

\bibitem{r27}
Cui J, Li X, Shao Z, et al. A Survey on Legal Large Language Models. arXiv preprint. 2024. arXiv:2312.03718, doi: 10.48550/arXiv.2312.03718.

\bibitem{r28}
Cormack GV, Clarke CLA, Buettcher S. Reciprocal Rank Fusion Outperforms Condorcet and Individual Rank Learning Methods. In: Proceedings of the 32nd International ACM SIGIR Conference on Research and Development in Information Retrieval (SIGIR 2009); 2009. p. 758-9, doi: 10.1145/1571941.1572114.

\end{thebibliography}

\end{document}
